\chapter{Timothy}
Timothy, the tiger cub, was discovered by Grandfather on a hunting expedition in the Terai jungle near Dehra.
\par
Grandfather was no shikari, but as he knew the forests of the Siwalik Hills better than most people, he was presuaded to accompany the party-it consisted of Several Very Important Persons from to Delhi-to advise on the terrrain and the direction the beaters should take once a tiger had been spotted.
\par
The camp itself was sumptuous-seven large tents (one for each shikari), and dinning tent and a number of servants' tents. The dinner was very good, as Grandfather admitted afterwards; it was not often that one saw hot-water plates, finger glasses and seven or eight courses in a tent in the jungle! But that was how things were done in the days o the viceroys . . . There were also some fifteen elephants, four of them with howdahs for the shikaris, and the others specially trained for talking part in the beat.
\par
The sportsmen never saw a tiger, nor did they shoot anything else, though they saw a number of deer, peacocks and wild boars. They were giving up all hope of finding a tiger and were beginning to shoot at jackals, when Grandfather, strolling down the forest path at some distance from the rest of the party, discovered a little tiger about eighteen inches long, hiding among the intricate roots of a banyan tree. Grandfather picked up him and brought him home after the camp had broken up. He had the distinction of being the only member of the party to have bagged any game, dead or alive.
\par
At first, the tiger cub, who was named Timothy by Grandmother, was brought up entirely on milk given to him in a feeding bottle by our cook, Mahmoud. But the milk proved too rich for him, and he was put on a diet of raw mutton and cod liver oil, to be followed later by a more tempting diet of pigeons and rabbits.
\par
Timothy was provided with two companions-Todo the monkey, who was bold enough to pull the youg tiger by the tail and then climb up the curtains if Timothy lost his temper, and a small morgel puppy, found on the road by Grandfather, Initially, Timothy appeared to be quite afraid of the puppy and darted back with a spring if it came too near.
\par
He would make absurd dashes at it with his large forepaws and then retreat to a ridiculously safe distance.
\par
Finally, he allowed the puppy to crawl on his back and rest there! One of Timothy's favourite amusements was to stalk anyone who would play wit hhim, and so, when I came to live with Grandfather, I became one of the tiger's favourites. With a crafty looking in his glittering eyes, and his body crouching, he would creep closer and closer to me, suddenly making a dash for my feet, rolling over on his back and kicking with delight, and pretending to bite my ankles.
\par
He was by this time the size of a full-grown retriever, and when I took him our for walks, people on the road would give us a wide berth. When he pulled hard on his chain, I had dificulty keeping up with him. His favourite place in the house was the drawing room, and he would make himself comfortable on the long sofa, reclining there with great dignity, and snarling at anybody who tried to get him off.
\par
Timothy had clean habits, and would scrub his face with his paws exactly like a cat. He slept at night in the cook's quarters, and was always delighted at being let out by him in the morning.
\par
'One of these days,' declared Grandmother in her prophetic manner, 'we are going to find Timothy sitting on Mahmoud's bed, and no sign of the cook except his clothes and shoes!'
\par
Of course, I never came to that, but when Timothy was about six months old, a change came over him; he grew steadily less friendly. When out for a walk with me, he would try to steal away to stalk a cat or someone's pet Pekinese. Sometimes at night we would hear frenzied clucking from the poultry house, and in the morning there would be feathers lying all over the veranda.
\par
Timothy had to be chained up more often. And, finally, when he began to stalk Mahmoud about the house with what looked like villainous intent, Grandfather decided it was time to transfer him to a zoo.
\par
The nearest zoo was at Lucknow, two hundred miles away. Reserving a first-class compartment for himself and Timothy-no one would share a compartment with them-Grandfather took him to Lucknow, where the zoo authorities were only too glad to receive as a gift a well fed and fairly civilized tiger.



